%=============================================================================!
% 6.3  THE EULER-LAGRANGE EQUATION                                            !
%=============================================================================!
\section{The Euler--Lagrange equation}
\label{sec:la-euler}

Hamilton's principle picks out the true path among all paths with the
same ends. This section turns it into an equation that the true path
obeys at every instant, using two tools, and shows that for a body moving
in a potential the equation is Newton's second law.

\begin{toolbox}[Integration by parts]
The product rule, $(fg)' = f'g + fg'$, integrated from $a$ to $b$ by the
fundamental theorem of calculus (\cref{sec:mo-integrating}), gives
$[fg]_a^b = \int_a^b f'g\,\dd t + \int_a^b fg'\,\dd t$, where $[fg]_a^b =
f(b)g(b) - f(a)g(a)$, so
\begin{equation*}
   \int_a^b f\,g'\,\dd t = \bigl[fg\bigr]_a^b - \int_a^b f'\,g\,\dd t :
\end{equation*}
an integral of $f$ times the derivative of $g$ becomes one of the
derivative of $f$ times $g$, at the cost of the bracket. For example,
with $f = t$ and $g = \sin t$, so that $f' = 1$ and $g' = \cos t$,
\begin{equation*}
   \int_0^{\pi/2} t\cos t\,\dd t = \Bigl[t\sin t\Bigr]_0^{\pi/2}
   - \int_0^{\pi/2}\sin t\,\dd t = \frac{\pi}{2} - \Bigl[-\cos t
   \Bigr]_0^{\pi/2} = \frac{\pi}{2} - 1 ,
\end{equation*}
since $\sin\frac{\pi}{2} = 1$ and $\sin 0 = 0$, and $-\cos t$, whose
derivative is $\sin t$, rises from $-1$ to 0.
\end{toolbox}

\subsection*{The first-order change of the action}

Change the path $q(t)$ to $q(t) + \zeta\eta(t)$, with $\eta$ any smooth
function that is zero at both ends, $\eta(t_1) = \eta(t_2) = 0$, so that
the changed path keeps the same ends. At each instant the Lagrangian
changes from $\Lag(q, \dot{q})$ to $\Lag(q + \zeta\eta, \dot{q} +
\zeta\dot\eta)$.

\begin{toolbox}[The first-order change of a function of a path]
Read $\Lag(q + \zeta\eta, \dot{q} + \zeta\dot\eta)$, at a fixed instant,
as a function of $\zeta$ alone: as $\zeta$ grows from 0, the point $(q +
\zeta\eta, \dot q + \zeta\dot\eta)$ moves in a straight line in the
plane of the two variables of $\Lag$, changing at the rate $(\eta, \dot
\eta)$ per unit of $\zeta$: at a fixed instant $\eta$ and $\dot\eta$ are
two fixed numbers, and $\zeta$ plays the part that time played in
\cref{sec:en-gradient}.
The chain rule along a path of \cref{sec:en-gradient}, $\dd\Lag/\dd\zeta
= \grad\Lag\cdot(\eta, \dot\eta)$, gives at $\zeta = 0$
\begin{equation*}
   \frac{\dd\Lag}{\dd\zeta} = \frac{\partial\Lag}{\partial q}\,\eta
   + \frac{\partial\Lag}{\partial\dot{q}}\,\dot\eta ,
\end{equation*}
where $\partial\Lag/\partial\dot q$ is the partial derivative with respect
to the second variable, the first held fixed (\cref{sec:en-gradient}).
At a single instant $q$ and $\dot q$ are two numbers, and $\Lag$ is an
ordinary function of the two, like the height $h(x, y)$ of a map; holding
$q$ fixed here does not make $\dot q$ zero, because the link between them,
that $\dot q$ is the rate of change of $q$, appears only once a whole path
is put in. By
the Taylor series (\cref{sec:mo-taylor}), $\Lag$ at $\zeta$ is $\Lag$ at 0
plus $\zeta$ times this, plus a remainder no larger than a
constant times $\zeta^2$.
\end{toolbox}

Integrating over the path, these remainders add up to at most $(t_2 -
t_1)$ times the largest such constant times $\zeta^2$, still of order
$\zeta^2$, so the action changes by
\begin{equation*}
   \Act(\zeta) - \Act(0) = \zeta\int_{t_1}^{t_2}\Bigl(\frac{\partial\Lag}
   {\partial q}\,\eta + \frac{\partial\Lag}{\partial\dot{q}}\,\dot\eta
   \Bigr)\dd t + \order{\zeta^2} .
\end{equation*}

The integral multiplying $\zeta$ contains both $\eta$ and $\dot\eta$. Along
the path, $\partial\Lag/\partial\dot{q}$ takes one value at each instant,
found by putting that instant's $q(t)$ and $\dot{q}(t)$ into it, so it is
a function of $t$ alone; its rate of change with time is written
$\dd(\partial\Lag/\partial\dot{q})/\dd t$, with an ordinary $\dd$
because nothing is now held fixed (for $\Lag = \frac12 m\dot{x}^2 - U$
below it is $\dd(m\dot{x})/\dd t = m\ddot{x}$). Integrating the second
term by parts, with $f = \partial\Lag/\partial\dot{q}$ and $g = \eta$,
\begin{equation*}
   \int_{t_1}^{t_2}\frac{\partial\Lag}{\partial\dot{q}}\,\dot\eta\,\dd t
   = \Bigl[\frac{\partial\Lag}{\partial\dot{q}}\,\eta\Bigr]_{t_1}^{t_2}
   - \int_{t_1}^{t_2}\frac{\dd}{\dd t}\Bigl(\frac{\partial\Lag}
   {\partial\dot{q}}\Bigr)\,\eta\,\dd t ,
\end{equation*}
and the bracket is zero, because $\eta$ is zero at both ends. Putting what
is left in place of the second term, and gathering the two integrals into
one, since both contain the factor $\eta$,
\begin{equation*}
   \Act(\zeta) - \Act(0) = \zeta\int_{t_1}^{t_2}\Bigl[\frac{\partial\Lag}
   {\partial q} - \frac{\dd}{\dd t}\Bigl(\frac{\partial\Lag}
   {\partial\dot{q}}\Bigr)\Bigr]\eta\,\dd t + \order{\zeta^2} .
\end{equation*}

\subsection*{The equation}

For the action to be stationary, its change must be of order $\zeta^2$
(\cref{sec:la-action}), so the term in $\zeta$ must vanish for every bump
$\eta$. Call the square bracket $B(t)$. If $B$ were positive at
some instant $t^*$, it would be positive on a short interval around
$t^*$, since it changes smoothly; choose $\eta$ positive on that interval
and zero everywhere else, such as $\eta = \sin^2 u$ with $u = \pi(t -
a)/(b - a)$ between the ends $a$ and $b$ of the interval, whose slope
$(2\pi/(b - a))\sin u\cos u$ is zero at both ends, so that it joins the
zero outside smoothly; the integral of $B\eta$ is positive, not
zero. The same argument rules out a negative $B$. So $B(t) = 0$ at every
instant:
\begin{equation}
   \frac{\dd}{\dd t}\Bigl(\frac{\partial\Lag}{\partial\dot{q}}\Bigr)
   = \frac{\partial\Lag}{\partial q} .
   \label{eq:la-euler}
\end{equation}
This is the Euler--Lagrange equation. Conversely, on a path that obeys it
at every instant $B = 0$, so the term in $\zeta$ vanishes for every bump:
the equation and Hamilton's principle pick out the same paths. A system
with several generalised coordinates has one such equation for each, $\dd(\partial\Lag/\partial
\dot{q}_k)/\dd t = \partial\Lag/\partial q_k$, since each coordinate can
be given its own bump while the others are left alone.

\subsection*{Newton's law again}

For a body on a line in a potential energy $U(x)$, $\Lag = \frac12 m
\dot{x}^2 - U(x)$. The partial derivatives are $\partial\Lag/\partial
\dot{x} = m\dot{x}$, with $x$ held fixed, and $\partial\Lag/\partial x =
-U'(x)$, with $\dot{x}$ held fixed, so \cref{eq:la-euler} reads
\begin{equation*}
   \frac{\dd}{\dd t}\bigl(m\dot{x}\bigr) = -U'(x) ,
   \qquad\text{that is}\qquad m\ddot{x} = F(x) ,
\end{equation*}
Newton's second law with the force $F = -U'$ of \cref{sec:en-potential}.
For the thrown ball, $U = mgx$, so $B = -mg - m\ddot{x}$, which is zero on
the true path, where $\ddot{x} = -g$. The middle integral of
\cref{sec:la-action} is the term in $\zeta$ above, which by parts is $\int
B\eta\,\dd t$, and this is why it vanished. For $N$ atoms with
$\Lag = \sum_i\frac12 m_i\lvert\dot{\vb{r}}_i\rvert^2 - U(\vb{r}^N)$, the
$3N$ equations are $m_i\ddot{\vb{r}}_i = -\grad_iU$, the equations of
motion of \cref{sec:en-atoms}. Hamilton's principle therefore contains
Newton's laws for every system whose forces come from a potential energy,
and from here on either can be used.

\Needspace{10\baselineskip}
\begin{bridge}
The Euler--Lagrange equation is the condition that a functional, a number
computed from a whole function, be stationary: the bracket $\partial\Lag/
\partial q - \dd(\partial\Lag/\partial\dot{q})/\dd t$ is the functional
derivative $\delta\Act/\delta q(t)$. It is the same operation as the
derivative $\delta E[n]/\delta n(\vb{r})$ of the density functional with
respect to the electron density; setting it equal to a Lagrange
multiplier that holds the number of electrons fixed gives the Euler
equation of density-functional theory, $\delta E/\delta n(\vb{r}) = \mu$,
with the chemical potential $\mu$ as the
multiplier\cite{Parr1978} (\cref{sec:la-multipliers}).
\end{bridge}

\begin{keyidea}
Hamilton's principle holds for a path exactly when, at every instant,
$\dd(\partial\Lag/\partial\dot{q}_k)/\dd t = \partial\Lag/\partial q_k$
for each generalised coordinate. For forces from a potential this is
Newton's second law.
\end{keyidea}

\begin{selfcheck}
Write down the Euler--Lagrange equation for a block on a spring, $\Lag =
\frac12 m\dot{x}^2 - \frac12 kx^2$.
\end{selfcheck}
